\cotefile{162-5}
\scanpage{33}
\margin{notes du 10.9.68}

Tapis de Quillen.

I \quad \uline{Relations entre catégories et ensembles semi-simplic\deleted{\illegible{1}}a\deleted{\illegible{1}}ux}.

A toute catégorie $C$, on associe un ensemble semi-simpli\deleted{\unsure{v}}cial $S(C)$,
trouvant \deleted{\unsure{u}}a\deleted{n}insi un foncteur pleinement fidèle \deleted{de Cat}
\[ S;\ (\mathrm{Cat}) \longrightarrow (\mathrm{Ssimpl})\quad . \]
\deleted{Si C est associé à} Les systêmes locaux d'ensemble sur $SC$ correspondent
aux fancteurs sur $C$ qui transforment toute flêche en isomorphisme (i.e.
qu\deleted{\unsure{ç|t}}i se factorisent par le groupoïde associé à $C$). Les $H^i$ sur $SC$ d'un tel
systême local ($H^0$ pour ens, $H^1$ pour groupes, $H^i$ quelconques pour groupes
abéliens) s'interprètent en termes des foncteurs $\underleftarrow{\mathrm{lim}}^{(i)}$ déri\deleted{g}vés de $\underleftarrow{\mathrm{lim}}$ ,
ou si on préfère, des $H^i($ du \uline{topos} $C$ . On voit ainsi à quelle condition
un foncteur $C \to C'$ induit un\deleted{e} homotopisme $SC \to SC'$ : \deleted{i}en vertu du
critère cohomologique de Artin-Mazur, il f et s que pour tout systême
de coéfficients $F'$sur $C'$, l'homomorphisme naturel $\underset{C'}{\underleftarrow{\mathrm{lim}}}^{(i)}F' \longrightarrow \underset{C}{\underleftarrow{\mathrm{lim}}}^{(i\ )}F$
soit un isomorphisme (pour les i pour lesquels cela a un sens). \deleted{En termes}

A $C$ on peut asso\deleted{v}c\unsure{\deleted{i}e|\deleted{e}i}r le topos $\widetilde{C}$ , qui varie de façon \uline{covariante} avec $C$.
(NB le foncteur $C \mapsto \widetilde{C}$ n'a plus rien de pleinement fidèle, semble-t\deleted{\unsure{-}}il ??)
Les systêmes de coéfficients \inserted{ensemblistes} sur $C$ ( $\overset{\text{\inserted{déf}}}{=}$ les foncteurs \inserted{$C^o \longrightarrow \mathrm{Ens}$} transformant isomor-phismes en isomorphismes) correspondent aux faisceaux localement constant\deleted{,}s,
i.e. les objets localement constants de $\widetilde{C}$, définis intrinsèquement en termes
de $\widetilde{C}$ . Ainsi, le fait pour un foncteur $F:C \to C'$ d'induire un\deleted{e équivalence}
homotopisme $S(C) \to S(C')$ ne dépend que du morphisme de topos $\widetilde{F}: \widetilde{C} \to \widetilde{C'}$
induit, et signifie que pour tout faisceau localement constant $F'$ sur $C'$ i.
e. sur $\widetilde{C'}$ , les applications induites $H^i(\widetilde{C'},F') \longrightarrow H^i(\widetilde{C},\widetilde{F}^{*}(F'))$
sont des isomorphismes (pour les i pour lesquels cela a un\deleted{en} sens).

\margin{32}
